% =========================================================
\section{Boutons-poussoirs et commutateurs}

\subsection{Introduction}

Dans un système automatisé, l'opérateur doit pouvoir intervenir pour :

\begin{itemize}
    \item mettre l'installation en marche ;
    \item arrêter la machine ;
    \item choisir un mode de fonctionnement ;
    \item signaler un défaut ou déclencher un arrêt d'urgence.
\end{itemize}

Les \textbf{boutons-poussoirs} et les \textbf{commutateurs} sont les organes de commande manuelle les plus utilisés dans les pupitres et armoires électriques.

Ils appartiennent à la \textbf{chaîne d'information} et sont raccordés aux entrées de l'API.

\subsection{Le bouton-poussoir}

\subsubsection{Définition}

Un bouton-poussoir est un organe de commande manuelle qui modifie l'état d'un contact lorsqu'on appuie dessus.

Il revient généralement à sa position initiale dès que l'opérateur relâche le bouton.

\subsubsection{Symbole}

Le symbole normalisé d'un bouton-poussoir est le suivant :

\begin{itemize}
    \item bouton à contact \textbf{NO} : contact ouvert au repos, fermé lorsqu'on appuie ;
    \item bouton à contact \textbf{NC} : contact fermé au repos, ouvert lorsqu'on appuie.
\end{itemize}

\subsubsection{Types de boutons-poussoirs}

On distingue :

\begin{itemize}
    \item le bouton \textbf{à contact NO} (Normalement Ouvert) ;
    \item le bouton \textbf{à contact NC} (Normalement Fermé) ;
    \item le bouton \textbf{à contacts NO + NC} (double contact) ;
    \item le bouton \textbf{à verrouillage} (reste enfoncé) ;
    \item le bouton \textbf{à clé} (commande protégée).
\end{itemize}

\subsection{Bouton à contact NO}

\subsubsection{Principe}

Un bouton à contact NO est ouvert au repos. Lorsqu'on appuie dessus, le contact se ferme et laisse passer le signal.

\[
\text{Repos} \rightarrow \text{Contact ouvert}
\qquad
\text{Appui} \rightarrow \text{Contact fermé}
\]

\subsubsection{Utilisation}

Il est utilisé pour les commandes de type \textbf{Marche}, \textbf{Départ cycle} ou \textbf{Acquittement}.

\paragraph{Exemple :} bouton « Marche » sur un pupitre de commande.

\subsection{Bouton à contact NC}

\subsubsection{Principe}

Un bouton à contact NC est fermé au repos. Lorsqu'on appuie dessus, le contact s'ouvre et interrompt le signal.

\[
\text{Repos} \rightarrow \text{Contact fermé}
\qquad
\text{Appui} \rightarrow \text{Contact ouvert}
\]

\subsubsection{Utilisation}

Il est utilisé pour les commandes de type \textbf{Arrêt} ou \textbf{Arrêt d'urgence}.

\paragraph{Exemple :} bouton « Arrêt » sur un pupitre de commande.

\paragraph{Important :} en cas de rupture de fil, un contact NC ouvre le circuit et provoque l'arrêt. C'est un principe de sécurité fondamental.

\subsection{Bouton d'arrêt d'urgence}

\subsubsection{Définition}

Le bouton d'arrêt d'urgence est un bouton à contact NC, généralement à \textbf{verrouillage mécanique}. Il reste enfoncé après l'appui et doit être déverrouillé manuellement (par rotation ou traction) avant de pouvoir redémarrer.

\subsubsection{Couleur}

Selon la norme, il est :

\begin{itemize}
    \item de couleur \textbf{rouge} ;
    \item sur fond \textbf{jaune}.
\end{itemize}

\subsubsection{Rôle}

Il a pour fonction d'arrêter immédiatement et en toute sécurité l'installation en cas de danger.

\paragraph{Attention :} l'arrêt d'urgence doit agir directement sur la chaîne d'énergie ou via un relais de sécurité, et non uniquement par le programme de l'API.

\subsection{Les commutateurs et sélecteurs}

\subsubsection{Définition}

Un commutateur ou sélecteur possède plusieurs positions \textbf{stables}. Il conserve son état après l'action de l'opérateur, contrairement au bouton-poussoir qui revient à sa position initiale.

\subsubsection{Exemples d'utilisation}

\begin{itemize}
    \item sélecteur « Manuel / Automatique » ;
    \item sélecteur « Marche / Arrêt » ;
    \item sélecteur de vitesse ;
    \item commutateur de mode de cycle.
\end{itemize}

\paragraph{Exemple :} un sélecteur à deux positions qui reste en position « Automatique » jusqu'à ce que l'opérateur le tourne vers « Manuel ».

\subsection{Différence entre bouton-poussoir et commutateur}

\begin{table}[h]
\centering
\begin{tabularx}{\textwidth}{|X|X|}
\hline
\textbf{Bouton-poussoir} & \textbf{Commutateur / sélecteur} \\
\hline
Retour automatique au repos & Conserve sa position \\
\hline
Action momentanée & Action permanente \\
\hline
Utilisé pour Marche, Arrêt, AU & Utilisé pour choix de mode \\
\hline
Un ou deux contacts & Plusieurs positions possibles \\
\hline
\end{tabularx}
\caption{Comparaison bouton-poussoir / commutateur}
\end{table}

\subsection{Contacts NO et NC : rappel}

\begin{itemize}
    \item \textbf{NO} : ouvert au repos, fermé lorsqu'on actionne ;
    \item \textbf{NC} : fermé au repos, ouvert lorsqu'on actionne.
\end{itemize}

\paragraph{Règle importante :} pour la sécurité, on utilise toujours des contacts \textbf{NC} pour les arrêts et les arrêts d'urgence.

\subsection{Raccordement à l'API}

Les boutons et commutateurs sont raccordés aux entrées \textbf{TOR} de l'API.

\paragraph{Exemple :}

\begin{table}[h]
\centering
\begin{tabularx}{\textwidth}{|c|X|c|}
\hline
\textbf{Adresse} & \textbf{Organe} & \textbf{Type de contact} \\
\hline
I0.0 & Bouton Marche & NO \\
\hline
I0.1 & Bouton Arrêt & NC \\
\hline
I0.2 & Bouton Arrêt d'urgence & NC verrouillé \\
\hline
I0.3 & Sélecteur Manuel/Auto & NO \\
\hline
I0.4 & Sélecteur Auto & NO \\
\hline
\end{tabularx}
\caption{Exemple de raccordement de boutons et sélecteurs}
\end{table}

\subsection{Exemple d'application : démarrage d'un moteur}

On souhaite démarrer un moteur à l'aide de deux boutons :

\begin{itemize}
    \item un bouton \textbf{Marche} (NO) ;
    \item un bouton \textbf{Arrêt} (NC).
\end{itemize}

Le programme de l'API doit :

\begin{itemize}
    \item démarrer le moteur lorsque le bouton Marche est appuyé ;
    \item maintenir le moteur en marche après relâchement (auto-maintien logiciel) ;
    \item arrêter le moteur lorsque le bouton Arrêt est appuyé.
\end{itemize}

\[
\text{Marche (NO)} \rightarrow \text{API} \rightarrow \text{Moteur}
\]

\[
\text{Arrêt (NC)} \rightarrow \text{API} \rightarrow \text{Arrêt du moteur}
\]

\subsection{Bonnes pratiques}

\begin{itemize}
    \item Repérer clairement chaque bouton sur le pupitre ;
    \item Utiliser des couleurs normalisées (vert pour Marche, rouge pour Arrêt) ;
    \item Utiliser un contact NC pour tout arrêt ou arrêt d'urgence ;
    \item Vérifier le fonctionnement du bouton avant le raccordement ;
    \item Noter soigneusement les adresses dans le tableau des entrées ;
    \item Éviter de cumuler plusieurs fonctions sur un même bouton.
\end{itemize}

\subsection{Erreurs fréquentes à éviter}

\begin{itemize}
    \item Confondre NO et NC ;
    \item Utiliser un bouton NO pour un arrêt d'urgence ;
    \item Raccorder un bouton sur une entrée déjà utilisée ;
    \item Oublier de respecter la tension d'alimentation du bouton ;
    \item Ne pas tester le bouton après câblage ;
    \item Ne pas distinguer bouton momentané et bouton verrouillé.
\end{itemize}

\subsection{Conclusion}

Les boutons-poussoirs et les commutateurs sont les organes de commande manuelle de base dans un système automatisé.

Ils permettent à l'opérateur d'agir sur la partie commande : mise en marche, arrêt, choix de mode, acquittement, arrêt d'urgence.

Leur raccordement correct aux entrées de l'API, ainsi que le respect des règles de sécurité, sont essentiels pour la fiabilité et la sécurité de l'installation.

\subsubsection*{À retenir}

\begin{quote}
Un bouton-poussoir est un organe momentané, tandis qu'un commutateur conserve sa position. Les contacts NO sont utilisés pour les commandes, les contacts NC pour les arrêts et les sécurités.
\end{quote}

\paragraph{Schéma fondamental :}

\[
\text{Boutons / commutateurs} \rightarrow \text{Entrées API} \rightarrow \text{Programme}
\]